% ============================================
% Johns Hopkins Letter Template
% Assistant Professor Cover Letter – Cal Poly (Orfalea College of Business)
% ============================================

\documentclass[11pt,letterpaper]{letter}

\usepackage{graphicx}
\usepackage{eso-pic}
\usepackage{geometry}
\usepackage{microtype}
\usepackage[hidelinks]{hyperref}

% ----- Page Setup -----
\geometry{left=25mm,right=25mm,top=25mm,bottom=25mm}

% ----- Logo Placement (first page only) -----
\newcommand{\PutJHULogoFG}[1]{%
  \AddToShipoutPictureFG*{%
    \AtPageUpperLeft{%
      \hspace{10mm}%
      \raisebox{-38mm}{%
        \includegraphics[width=6cm]{#1}%
      }%
    }%
  }%
}

% ----- Sender Information -----
\signature{Decory Edwards}
\address{Department of Economics \\ Johns Hopkins University \\ Baltimore, MD 21218 \\ \texttt{dedwar65@jh.edu}}

\begin{document}
\PutJHULogoFG{Images/jhulogo.png}

\begin{letter}{Search Committee \\
Orfalea College of Business \\
California Polytechnic State University, San Luis Obispo \\
San Luis Obispo, CA 93407 \\ USA}

\opening{Dear Members of the Search Committee,}

I am writing to express my interest in the tenure-track Assistant Professor position in Economics at the Orfalea College of Business at Cal Poly, beginning August 17, 2026. I am a Ph.D. candidate in Economics at Johns Hopkins University, advised by Professors Christopher D.~Carroll, Nicholas W.~Papageorge, and Stelios Fourakis, and I expect to receive my degree in May 2026. My research fields are macroeconomics, wealth and income dynamics, and computational economics.

My research agenda focuses on understanding the determinants of wealth inequality and the mechanisms through which heterogeneity in returns to wealth shapes aggregate outcomes. In my job-market paper, I estimate the degree of heterogeneity in individual portfolio returns using micro-data from the Survey of Consumer Finances and simulate a structural model in which households face idiosyncratic return risk. I show that allowing for persistent heterogeneity in returns substantially improves the model’s ability to match the empirical distribution of wealth, offering a tractable way to connect micro-level return dispersion to macroeconomic inequality.

In a second project, I use the Health and Retirement Study to examine the relationship between interpersonal trust and portfolio performance. Building on the literature that finds a hump-shaped relationship between trust and income, I show that a similar relationship holds for individual returns to wealth measured in the HRS data, highlighting a behavioral channel through which social attitudes shape saving and investment outcomes. This work also provides new estimates of the persistent component of returns to net worth, which motivate the calibration of heterogeneity in my structural model.

My work in quantitative macroeconomics, heterogeneous-agent modeling, and empirical analysis of wealth and portfolio outcomes fits naturally with these priorities and with teaching undergraduate and master’s-level courses that build strong empirical and analytical skills. I am especially drawn to Cal Poly’s commitment to inclusive, accessible instruction and to providing students with practical, data-driven tools that prepare them for careers in business, policy, and further study.

\textbf{Teaching.} My teaching experience centers on undergraduate macroeconomics and an upper-level macro policy seminar, where I have led weekly discussion sections and held regular office hours for classes ranging from 16 to 40 students. My approach is student-centered: I aim to help students “convince themselves” that they have moved from confusion to understanding by holding structured office-hour coaching that builds trust and confidence. At Cal Poly, I am prepared to teach \textit{Principles of Macroeconomics}, \textit{Principles of Microeconomics}, \textit{Econometrics}, and \textit{Applied Econometrics} at both the undergraduate and master’s levels, along with field courses in applied macroeconomics and applied microeconomics as needed.

I believe I would be a strong fit because my computational and empirical toolkit allows me to (i) produce rigorous, data-backed estimates of returns and distributional outcomes, (ii) build and validate quantitative models that speak to macroeconomic and business-relevant questions, and (iii) mentor students effectively in a setting that values inclusive teaching, in-person engagement, and active scholarship. I am enthusiastic about the opportunity to contribute to Cal Poly’s teaching, research, and service missions in the Orfalea College of Business.

I would be honored to join the economists in the business school at Cal Poly. Thank you for your consideration; I welcome the opportunity to discuss my fit with your department at your convenience.

\closing{Sincerely,}

\end{letter}
\end{document}
